% c19 原书 p.236–246（PDF p.243–253）
\begin{lemma}\label{it:2.7.14}
设 $u\colon\fF\to\fG$ 是 $\cO_{\widehat{X}}$-模的态射。假设 $\fG$、$\ker u$ 与 $\coker u$ 都是可代数化的凝聚 $\cO_{\widehat{X}}$-模。则 $\fF$ 是可代数化的凝聚 $\cO_{\widehat{X}}$-模。
\end{lemma}

\begin{proof}
对态射 $\fG\to\coker u$ 应用 2.7.12，我们看到 $\Image u$ 是凝聚且可代数化的。对短正合序列
\begin{equation*}
0\to\ker u\to\fF\to\Image u\to 0
\end{equation*}
应用 2.7.13，我们看到 $\fF$ 是凝聚且可代数化的。
\end{proof}

\begin{lemma}\label{it:2.7.15}
考虑诺特概形态射的交换图
\begin{center}
\begin{tikzcd}
Y \arrow[r,"g"] \arrow[dr] & Z \arrow[d] \\
& \Spec A.
\end{tikzcd}
\end{center}
设 $\widehat{Y}$ 与 $\widehat{Z}$ 分别是 $Y$ 与 $Z$ 关于理想层 $I^{\sim}\cO_Y$ 与 $I^{\sim}\cO_Z$ 的形式完全化，$\hat{g}\colon\widehat{Y}\to\widehat{Z}$ 是由 $g$ 诱导的态射。假设 $g$ 是固有态射。若 $\fF$ 是可代数化的凝聚 $\cO_{\widehat{Y}}$-模，则 $\hat{g}_*\fF$ 是可代数化的凝聚 $\cO_{\widehat{Z}}$-模。
\end{lemma}

\begin{proof}
设 $\cF$ 是凝聚 $\cO_Y$-模，使得 $\fF\cong\widehat{\cF}$。由 2.7.6 我们有
\begin{equation*}
\hat{g}_*\fF\cong\hat{g}_*\widehat{\cF}\cong\widehat{g_*\cF}.
\end{equation*}
由 2.5.1，$g_*\cF$ 是凝聚 $\cO_Z$-模。所以 $\hat{g}_*\fF$ 是凝聚且可代数化的。
\end{proof}

\begin{lemma}\label{it:2.7.16}
设 $g\colon Y\to Z$ 是诺特概形的固有态射，$\cI$ 是 $\cO_Z$ 的凝聚理想层，$\widehat{Y}$ 与 $\widehat{Z}$ 分别是 $Y$ 与 $Z$ 关于理想层 $\cI\cO_Y$ 与 $\cI$ 的形式完全化，$\hat{g}\colon\widehat{Y}\to\widehat{Z}$ 是由 $g$ 诱导的态射。假设 $U$ 是 $Z$ 的开子概形，且 $g$ 诱导了 $g^{-1}(U)$ 与 $U$ 的同构。设 $\cJ$ 是 $\cO_Z$ 的凝聚理想层，使得 $Z-U=\Supp(\cO_Z/\cJ)$。则对每个凝聚 $\cO_{\widehat{Z}}$-模 $\fG$，标准态射
\begin{equation*}
\fG\to\hat{g}_*\hat{g}^*\fG
\end{equation*}
的核与余核都被 $\widehat{\cJ}^n$ 零化，对某个自然数 $n$。
\end{lemma}

\begin{proof}
问题关于 $Z$ 是局部的。我们可以假设 $Z=\Spec B$，其中 $B$ 是诺特环。设 $I$ 与 $J$ 是 $B$ 的分别对应于 $\cI$ 与 $\cJ$ 的理想。设 $\widehat{B}$ 是 $B$ 的 $I$-进完全化，$J'=J\widehat{B}$，$Z'=\Spec\widehat{B}$，$U'$ 是 $U$ 在 $Z'$ 中的原像，$Y'=Y\times_Z Z'$，$g'\colon Y\times_Z Z'\to Z'$ 是投影。
\begin{center}
\begin{tikzcd}
\widehat{Y} \arrow[r,"\hat{g}"] \arrow[d] & \widehat{Z} \arrow[d] \\
Y'=Y\times_Z Z' \arrow[r,"g'"] \arrow[d] & Z' \arrow[d] \\
Y \arrow[r,"g"] & Z.
\end{tikzcd}
\end{center}
则 $g'$ 诱导了 $g'^{-1}(U')$ 与 $U'$ 的同构，且 $Z'-U'=\Supp(\cO_{Z'}/J'^{\sim})$。注意 $\widehat{Y}$ 与 $\widehat{Z}$ 也是 $Y'$ 与 $Z'$ 关于理想层 $(\widehat{I})^{\sim}$ 与 $(\widehat{I})^{\sim}\cO_{Y'}$ 的形式完全化。此外，由 1.5.14 (iv)，对每个 $n$ 我们有
\begin{equation*}
\widehat{(J'^n)^{\sim}}\cong(J'^n)^{\sim}\cO_{\widehat{Z}}=(J^n)^{\sim}\cO_{\widehat{Z}}\cong\widehat{(J^n)^{\sim}},
\end{equation*}
即
\begin{equation*}
\widehat{(J'^n)^{\sim}}\cong\widehat{(J^n)^{\sim}}.
\end{equation*}
要证明引理，我们可以因此把 $Z$ 换成 $Z'$。于是我们可以假设 $B$ 是 $I$-进完备的。则由 1.5.20，存在有限生成 $B$-模 $M$，使得 $\fG\cong(M^{\sim})^{\wedge}$。设 $i_Y\colon\widehat{Y}\to Y$ 与 $i_Z\colon\widehat{Z}\to Z$ 是标准态射。我们有
\begin{equation*}
\hat{g}^*\fG\cong\hat{g}^*((M^{\sim})^{\wedge})\cong\hat{g}^*i_Z^*M^{\sim}\cong i_Y^*g^*M^{\sim}\cong(g^*M^{\sim})^{\wedge}.
\end{equation*}
所以由 2.7.6 有
\begin{equation*}
\hat{g}_*\hat{g}^*\fG\cong\hat{g}_*(g^*M^{\sim})^{\wedge}\cong(g_*g^*M^{\sim})^{\wedge}.
\end{equation*}
可以验证，标准态射
\begin{equation*}
\fG\to\hat{g}_*\hat{g}^*\fG
\end{equation*}
是由标准态射
\begin{equation*}
M^{\sim}\to g_*g^*M^{\sim}
\end{equation*}
经取形式完全化诱导的。由 2.5.1，$g_*g^*M^{\sim}$ 是凝聚的。设 $K$ 是有限生成 $B$-模，使得
\begin{equation*}
\ker(M^{\sim}\to g_*g^*M^{\sim})\cong K^{\sim}.
\end{equation*}
由于 $g$ 诱导了 $g^{-1}(U)$ 与 $U$ 的同构，$K^{\sim}$ 支撑在 $U$ 之外。由 1.4.14，存在整数 $n$，使得 $(J^n)^{\sim}K^{\sim}=0$。由 1.5.13 (i)，$(K^{\sim})^{\wedge}$ 同构于 $\fG\to\hat{g}_*\hat{g}^*\fG$ 的核。所以 $\fG\to\hat{g}_*\hat{g}^*\fG$ 的核被 $\widehat{\cJ}^n$ 零化。类似地，$\fG\to\hat{g}_*\hat{g}^*\fG$ 的余核被 $\widehat{\cJ}^n$ 零化，对某个自然数 $n$。
\end{proof}

\begin{lemma}\label{it:2.7.17}
设 $A$ 是诺特环，$I$ 是 $A$ 的理想，$S=\Spec A$，$f\colon X\to S$ 是射影态射，$\cO_X(1)$ 是 $S$ 上极丰沛的可逆 $\cO_X$-模，$\widehat{X}$ 是 $X$ 关于理想层 $I^{\sim}\cO_X$ 的形式完全化，$\fF$ 是凝聚 $\cO_{\widehat{X}}$-模。对每个 $k$，设 $X_k$ 是 $X$ 的由理想层 $(I^k)^{\sim}\cO_X$ 定义的闭子概形，$i_k\colon X_k\to X$、$\hat{i}_k\colon X_k\to\widehat{X}$ 与 $i\colon\widehat{X}\to X$ 是标准态射。令
\begin{align*}
\cO_X(n)&=\cO_X(1)^{\otimes n},\\
\cF_k&=\hat{i}_k^*\fF,\\
\cF_k(n)&=\cF_k\otimes_{\cO_{X_k}}i_k^*\cO_X(n),\\
\fF(n)&=\fF\otimes_{\cO_{\widehat{X}}}i^*\cO_X(n).
\end{align*}
则存在整数 $N$，使得对每个 $n\ge N$ 与每个 $k$，标准同态
\begin{equation*}
H^0(\widehat{X},\fF(n))\to H^0(X_k,\cF_k(n))
\end{equation*}
是满射，且对每个 $p\ge1$ 有
\begin{equation*}
H^p(\widehat{X},\fF(n))=0.
\end{equation*}
\end{lemma}

\begin{proof}
设 $B=\bigoplus_{k=0}^{\infty}I^k/I^{k+1}$。用如下笛卡尔图固定记号：
\begin{center}
\begin{tikzcd}
X'=X\times_S\Spec B \arrow[r,"g'"] \arrow[d,"f'"'] & X_1 \arrow[d,"f_1"] \\
\Spec B \arrow[r,"g"'] & \Spec A/I,
\end{tikzcd}
\end{center}
其中 $X_1$ 是 $X$ 的由理想层 $I^{\sim}\cO_X$ 定义的闭子概形，$f_1\colon X_1\to\Spec A/I$ 是由 $f$ 诱导的态射，$g$ 是由同态
\begin{equation*}
A/I\to B=A/I\oplus I/I^2\oplus\cdots
\end{equation*}
诱导的态射。注意 $g$ 与 $g'$ 是仿射态射，且
\begin{equation*}
g'_*\cO_{X'}\cong f_1^*B^{\sim}\cong\bigoplus_{k=0}^{\infty}f_1^*(I^k/I^{k+1})^{\sim}.
\end{equation*}
考虑 $\widehat{X}$ 上的层 $\bigoplus_{k=0}^{\infty}I^k\fF/I^{k+1}\fF$。由于每个 $I^k\fF/I^{k+1}\fF$ 被 $I$ 零化且 $\cO_{X_1}\cong\cO_{\widehat{X}}/I\cO_{\widehat{X}}$，我们可以把 $I^k\fF/I^{k+1}\fF$ 看作 $\cO_{X_1}$-模。它是凝聚 $\cO_{X_1}$-模。事实上，我们有 $\cO_{\widehat{X}}$-模的正合序列
\begin{equation*}
0\to I^k\fF/I^{k+1}\fF\to\fF/I^{k+1}\fF\to\fF/I^k\fF\to 0.
\end{equation*}
由于 $\cO_{X_{k+1}}\cong\cO_{\widehat{X}}/I^{k+1}\cO_{\widehat{X}}$ 且序列中每个层都被 $I^{k+1}$ 零化，我们可以把该序列看作 $\cO_{X_{k+1}}$-模的正合序列。此时 $\fF/I^{k+1}\fF$ 与 $\cF_{k+1}$ 等同，$\fF/I^k\fF$ 与 $(i_{k,k+1})_*\cF_k$ 等同，其中 $i_{k,k+1}\colon X_k\to X_{k+1}$ 是标准态射。由 1.5.17，$\cF_k$ 与 $\cF_{k+1}$ 是凝聚的。所以 $I^k\fF/I^{k+1}\fF$ 是凝聚 $\cO_{X_{k+1}}$-模。这说明它是凝聚 $\cO_{X_1}$-模。于是 $\bigoplus_{k=0}^{\infty}I^k\fF/I^{k+1}\fF$ 是拟凝聚 $\cO_{X_1}$-模。它也是 $\bigoplus_{k=0}^{\infty}f_1^*(I^k/I^{k+1})^{\sim}$-模。于是由 1.4.12，存在拟凝聚 $\cO_{X'}$-模 $\fF'$，使得
\begin{equation*}
g'_*\fF'\cong\bigoplus_{k=0}^{\infty}I^k\fF/I^{k+1}\fF.
\end{equation*}
注意 $\fF'$ 实际上是凝聚的。对每个 $n$，由 1.4.26 我们有
\begin{equation*}
g'_*(\fF'(n))\cong\bigoplus_{k=0}^{\infty}I^k\fF(n)/I^{k+1}\fF(n),
\end{equation*}
其中 $\fF'(n)$ 是 $\fF'$ 被 $\cO_X(n)$ 在 $X'$ 中的原像扭得到的层。由于 $g'$ 是仿射的，双正则谱序列
\begin{equation*}
H^p(X_1,R^qg'_*(\fF'(n)))\Rightarrow H^{p+q}(X',\fF'(n))
\end{equation*}
由 2.4.8 退化，从而
\begin{equation*}
H^p(X_1,g'_*(\fF'(n)))\cong H^p(X',\fF'(n)),
\end{equation*}
即
\begin{equation*}
\bigoplus_{k=0}^{\infty}H^p(X_1,I^k\fF(n)/I^{k+1}\fF(n))\cong H^p(X',\fF'(n)).
\end{equation*}
由 2.4.12 (ii)，存在整数 $N$，使得对每个 $n\ge N$ 与每个 $p\ge1$ 有 $H^p(X',\fF'(n))=0$。所以我们有
\begin{equation*}
H^p(X_1,I^k\fF(n)/I^{k+1}\fF(n))=0
\end{equation*}
对每个 $k\ge0$、$p\ge1$ 与 $n\ge N$。对 $k$ 归纳，可以证明：若 $k\ge1$、$p\ge1$ 且 $n\ge N$，则
\begin{equation*}
H^p(\widehat{X},\fF(n)/I^k\fF(n))=0,
\end{equation*}
即
\begin{equation*}
H^p(X_k,\cF_k(n))=0.
\end{equation*}
特别地，逆向系统 $(H^p(X_k,\cF_k(n)))_{k\in\bN}$ 对每个 $p\ge1$ 与每个 $n\ge N$ 都满足 Mittag-Leffler 条件。另一方面，与层的短正合序列
\begin{equation*}
0\to I^k\fF(n)/I^{k+1}\fF(n)\to\fF(n)/I^{k+1}\fF(n)\to\fF(n)/I^k\fF(n)\to 0
\end{equation*}
相伴，我们有上同调的正合序列
\begin{equation*}
H^0(X_{k+1},\cF_{k+1}(n))\to H^0(X_k,\cF_k(n))\to H^1(X_1,I^k\fF(n)/I^{k+1}\fF(n)).
\end{equation*}
由于对每个 $k\ge0$ 与每个 $n\ge N$ 有 $H^1(X_1,I^k\fF(n)/I^{k+1}\fF(n))=0$，同态
\begin{equation*}
H^0(X_{k+1},\cF_{k+1}(n))\to H^0(X_k,\cF_k(n))
\end{equation*}
是满射，从而逆向系统 $(H^0(X_k,\cF_k(n)))_{k\in\bN}$ 对每个 $n\ge N$ 也满足 Mittag-Leffler 条件。此外，对 $X_1$ 的每个仿射开子概形 $U$ 与每对 $l\ge k$，同态
\begin{equation*}
H^0(U,\cF_l(n))\to H^0(U,\cF_k(n))
\end{equation*}
是满射，且对每个 $p\ge1$ 有
\begin{equation*}
H^p(U,\cF_k(n))=0.
\end{equation*}
于是由 2.7.3 有
\begin{equation*}
H^p(\widehat{X},\fF(n))\cong\plim_k H^p(X_k,\cF_k(n))
\end{equation*}
对每个 $p$ 与每个 $n\ge N$。由这个公式与逆向系统 $(H^p(X_k,\cF_k(n)))_{k\in\bN}$ 的性质，我们看到对每个 $n\ge N$ 与每个 $k$，同态
\begin{equation*}
H^0(\widehat{X},\fF(n))\to H^0(X_k,\cF_k(n))
\end{equation*}
是满射，且对每个 $p\ge1$ 有
\begin{equation*}
H^p(\widehat{X},\fF(n))=0.
\end{equation*}
\end{proof}

\begin{corollary}\label{it:2.7.18}
保持 2.7.17 的假设与记号。存在整数 $N$，使得对每个 $n\ge N$，$\fF(n)$ 由有限多个整体截面生成，即存在某个自然数 $m$ 的满射 $\cO_{\widehat{X}}^m\to\fF(n)$。
\end{corollary}

\begin{proof}
由 1.4.25，存在整数 $N$，使得对每个 $n\ge N$，$\cF_1(n)$ 由有限多个整体截面生成。此外，由 2.7.17，我们可以取 $N$ 足够大，使得
\begin{equation*}
\Gamma(\widehat{X},\fF(n))\to\Gamma(X_1,\cF_1(n))
\end{equation*}
是满射。设 $s_1,\dots,s_m\in\Gamma(\widehat{X},\fF(n))$，使得它们在 $\Gamma(X_1,\cF_1(n))$ 中的像生成 $\cF_1(n)$。考虑态射
\begin{equation*}
\cO_{\widehat{X}}^m\to\fF(n),\quad e_i\mapsto s_i,
\end{equation*}
其中对每个 $i$，$e_i$ 是 $\cO_{\widehat{X}}^m$ 的第 $i$ 个分量为 $1$、其余分量为 $0$ 的截面。我们来证明 $\cO_{\widehat{X}}^m\to\fF(n)$ 是满射。对 $X$ 的每个仿射开子概形 $U=\Spec B$，设 $\widehat{B}$ 是 $B$ 的 $(IB)$-进完全化。用 $\widehat{B}$ 替换 $B$ 不会影响我们的问题。于是我们可以假设 $B$ 是 $(IB)$-进完备的。由 1.5.20，我们有 $\fF(n)|_{U\cap\widehat{X}}\cong(M^{\sim})^{\wedge}$，其中 $M$ 是某个有限生成 $B$-模。同态
\begin{equation*}
\cO_{\widehat{X}}^m(U\cap\widehat{X})\to\fF(n)(U\cap\widehat{X})
\end{equation*}
可以等同于同态
\begin{equation*}
B^m\to M.
\end{equation*}
由 $\cO_{\widehat{X}}^m\to\fF(n)$ 的选取，同态
\begin{equation*}
(B/IB)^m\to M/IM
\end{equation*}
是满射。由 1.5.7 有 $IB\subset\operatorname{rad}(B)$。于是由中山引理，$B^m\to M$ 是满射。因此
\begin{equation*}
((B^m)^{\sim})^{\wedge}\to(M^{\sim})^{\wedge}
\end{equation*}
由 1.4.2 (iii) 与 1.5.13 (i) 是满射，即 $\cO_{\widehat{X}}^m\to\fF(n)$ 在 $U\cap\widehat{X}$ 上的限制是满射。由于 $U$ 任意，$\cO_{\widehat{X}}^m\to\fF(n)$ 是满射。
\end{proof}

\begin{tnote}
原书 p.241 把 $\Gamma(X_1,\cF_1(n))$ 印作 $\Gamma(X,\cF_1(n))$，落了下标 $1$。这里只能取 $X_1$：$\cF_1$ 是 $X_1$ 上的模，且原书紧接的下一句自己就写「它们在 $\Gamma(X,\cF_1(n))$ 中的像」，而 2.7.17 给出的正是 $k=1$ 时的满射。译文按 $X_1$ 处理。[依据：原书 p.241 第 2 段；参 2.7.17]
\end{tnote}

\begin{corollary}\label{it:2.7.19}
若 $X$ 在 $\Spec A$ 上射影，则 2.7.10 成立。
\end{corollary}

\begin{proof}
固定一个在 $\Spec A$ 上极丰沛的可逆 $\cO_X$-模 $\cO_X(1)$。对每个凝聚 $\cO_{\widehat{X}}$-模 $\fF$，由 2.7.18 我们可以找到满射 $\cO_{\widehat{X}}(n)^m\to\fF$。对该满射的核再次应用 2.7.18，我们看到存在如下形式的正合序列
\begin{equation*}
\cO_{\widehat{X}}(n')^{m'}\to\cO_{\widehat{X}}(n)^m\to\fF\to 0
\end{equation*}
（对某些自然数 $m$、$n$、$m'$、$n'$）。$\cO_{\widehat{X}}(n)^m$ 与 $\cO_{\widehat{X}}(n')^{m'}$ 都是可代数化凝聚 $\cO_{\widehat{X}}$-模，因为它们分别是 $\cO_X(n)^m$ 与 $\cO_X(n')^{m'}$ 的形式完全化。由 2.7.12，$\fF$ 也可代数化。
\end{proof}

\begin{tnote}
2.7.12–2.7.19 是为定理 2.7.10 的存在性准备的下降链条：2.7.12、2.7.13 说明「可代数化」对核、余核与扩张封闭；2.7.14 补上中间项这一情形；2.7.15、2.7.16 允许沿固有态射（由 Chow 引理给出的射影覆盖）下降；2.7.17–2.7.19 则在射影情形下给出足够的整体截面，从而完成可代数化。[依据：原书 p.235–242]
\end{tnote}

最后，我们来证明 2.7.10。

\begin{proof}[定理 2.7.10 的证明]
设 $\mathcal{S}$ 是满足下述条件的 $X$ 的闭子集 $Z$ 的族：存在 $X$ 的闭子概形 $(Z,\cO_Z)$，其底拓扑空间为 $Z$，且存在不可代数化的凝聚 $\cO_Z$-模。若 $\mathcal{S}=\varnothing$，则我们完成。否则，由于 $X$ 是诺特的，$\mathcal{S}$ 有极小元。设 $Z$ 是 $X$ 的闭子概形，使得存在不可代数化的凝聚 $\cO_Z$-模 $\fF$，且 $Z$ 的底拓扑空间是 $\mathcal{S}$ 的极小元。对 $Z\to\Spec A$ 应用 Chow 引理 1.4.18，我们得到射影态射 $g\colon Y\to Z$，使得 $Y$ 在 $\Spec A$ 上射影，且存在 $Z$ 的稠密开子集 $U$，使得 $g$ 诱导了 $g^{-1}(U)$ 与 $U$ 的同构。由 1.5.22 (ii)，$\hat{g}^*\fF$ 是凝聚的。由于 $Y$ 在 $\Spec A$ 上射影，由 2.7.19，$\hat{g}^*\fF$ 可代数化。于是由 2.7.15，$\hat{g}_*\hat{g}^*\fF$ 可代数化。设 $\cJ$ 是 $\cO_Z$ 的凝聚理想层，使得 $Z-U=\Supp(\cO_Z/\cJ)$。设 $\mathcal{R}$ 是标准态射
\begin{equation*}
\fF\to\hat{g}_*\hat{g}^*\fF
\end{equation*}
的核。由 2.7.16，$\mathcal{R}$ 被 $\widehat{\cJ}^n$ 零化，对某个自然数 $n$。设
\begin{equation*}
h\colon(Z-U,(\cO_Z/\cJ^n)|_{Z-U})\to(Z,\cO_Z)
\end{equation*}
是闭浸入。我们断言 $\mathcal{R}=\hat{h}_*\hat{h}^*\mathcal{R}$。事实上，由 1.5.13 (i)，序列
\begin{equation*}
0\to\widehat{\cJ}^n\to\cO_{\widehat{Z}}\to(h_*\cO_{Z-U})^{\wedge}\to 0
\end{equation*}
正合。由 2.7.6（或直接证明），我们有
\begin{equation*}
(h_*\cO_{Z-U})^{\wedge}\cong\hat{h}_*\cO_{\widehat{Z-U}},
\end{equation*}
其中 $\cO_{\widehat{Z-U}}$ 是 $\cO_{Z-U}=(\cO_Z/\cJ^n)|_{Z-U}$ 的形式完全化。所以我们有正合序列
\begin{equation*}
0\to\widehat{\cJ}^n\to\cO_{\widehat{Z}}\to\hat{h}_*\cO_{\widehat{Z-U}}\to 0.
\end{equation*}
所以
\begin{equation*}
\cO_{\widehat{Z-U}}\cong(\cO_{\widehat{Z}}/\widehat{\cJ}^n)|_{\widehat{Z-U}}.
\end{equation*}
此外，由于 $\mathcal{R}$ 被 $\widehat{\cJ}^n$ 零化，它支撑在 $\widehat{Z-U}$ 上。（事实上，若 $V$ 是 $\widehat{Z}$ 的与 $\widehat{Z-U}$ 不交的开子集，则 $1$ 是 $\widehat{\cJ}^n(V)$ 中的截面。但对 $\mathcal{R}(V)$ 中的任一截面 $x$，我们有 $x=1\cdot x=0$。）所以我们有
\begin{equation*}
\mathcal{R}\cong\hat{h}_*\hat{h}^*\mathcal{R}.
\end{equation*}
由 1.5.22 (ii)，$\hat{h}^*\mathcal{R}$ 是凝聚的。由 $Z$ 在 $\mathcal{S}$ 中的底拓扑空间的极小性，$\hat{h}^*\mathcal{R}$ 可代数化。于是由 2.7.15（或直接证明），$\mathcal{R}\cong\hat{h}_*\hat{h}^*\mathcal{R}$ 可代数化，即态射 $\fF\to\hat{g}_*\hat{g}^*\fF$ 的核可代数化。类似地，其余核也可代数化。对态射 $\fF\to\hat{g}_*\hat{g}^*\fF$ 应用 2.7.14，我们看到 $\fF$ 可代数化。这与 $\fF$ 被选为不可代数化矛盾。所以 $\mathcal{S}=\varnothing$。这就证明了我们的定理。
\end{proof}

\begin{tnote}
定理 2.7.10 的证明取 Noether 归纳：以「存在不可代数化的凝聚层」的闭子集族 $\mathcal{S}$ 的极小元为反例，再用 Chow 引理把它换成射影覆盖上已经可代数化的情形，最后用 2.7.14 反推矛盾。这是 Grothendieck 存在定理证明的经典结构。[依据：原书 p.242–243]
\end{tnote}

\begin{theorem}\label{it:2.7.20}
设 $A$ 是诺特环，$I$ 是 $A$ 的理想，$S=\Spec A$，$X$ 与 $Y$ 是 $S$ 上的两个固有概形。对每个 $k$，设 $S_k=\Spec A/I^k$，$X_k=X\times_S S_k$，$Y_k=Y\times_S S_k$。若 $A$ 是 $I$-进完备的，则我们有标准双射
\begin{equation*}
\Hom_S(X,Y)\cong\plim_k\Hom_{S_k}(X_k,Y_k),
\end{equation*}
其中 $\Hom_S(X,Y)$ 是从 $X$ 到 $Y$ 的 $S$-态射的集合，$\Hom_{S_k}(X_k,Y_k)$ 是从 $X_k$ 到 $Y_k$ 的 $S_k$-态射的集合。
\end{theorem}

\begin{proof}
对每个 $S$-态射 $f\colon X\to Y$ 与每个 $k$，设 $f_k\colon X_k\to Y_k$ 是由 $f$ 诱导的 $S_k$-态射。则 $(f_k)\in\plim_k\Hom_{S_k}(X_k,Y_k)$。我们证明映射
\begin{equation*}
\Hom_S(X,Y)\to\plim_k\Hom_{S_k}(X_k,Y_k),\quad f\mapsto(f_k)
\end{equation*}
是双射。设 $f,g\colon X\to Y$ 是使它们对每个 $k$ 诱导相同的态射 $X_k\to Y_k$ 的两个 $S$-态射。考虑它们的图像
\begin{equation*}
\Gamma_f,\Gamma_g\colon X\to X\times_S Y.
\end{equation*}
由于 $Y$ 在 $S$ 上分离，它们是闭浸入。设 $\sJ_f$ 与 $\sJ_g$ 是它们的理想层。则由 $f$ 与 $g$ 诱导的态射 $X_k\to Y_k$ 的图像理想层是
\begin{equation*}
\sJ_f\cO_{X_k\times_{S_k}Y_k}=\sJ_g\cO_{X_k\times_{S_k}Y_k}.
\end{equation*}
我们有
\begin{equation*}
\cO_{X_k\times_{S_k}Y_k}/\sJ_f\cO_{X_k\times_{S_k}Y_k}=\cO_{X_k\times_{S_k}Y_k}/\sJ_g\cO_{X_k\times_{S_k}Y_k}.
\end{equation*}
从而我们有同构的相容族
\begin{equation*}
i_k^*(\cO_{X\times_S Y}/\sJ_f)\cong i_k^*(\cO_{X\times_S Y}/\sJ_g),
\end{equation*}
使得图表
\begin{center}
\begin{tikzcd}[column sep=1.1em,row sep=1.7em]
\cO_{X_k\times_{S_k}Y_k}/\sJ_f\cO_{X_k\times_{S_k}Y_k} \arrow[dd,equal] \arrow[drr] & & & & i_k^*(\cO_{X\times_S Y}/\sJ_f) \arrow[dd,"\cong"] \arrow[dll] \\
& & \cO_{X_k\times_{S_k}Y_k}\cong i_k^*\cO_{X\times_S Y} & & \\
\cO_{X_k\times_{S_k}Y_k}/\sJ_g\cO_{X_k\times_{S_k}Y_k} \arrow[urr] & & & & i_k^*(\cO_{X\times_S Y}/\sJ_g) \arrow[ull]
\end{tikzcd}
\end{center}
交换，其中 $i_k\colon X_k\times_{S_k}Y_k\to X\times_S Y$ 是标准态射。由 1.5.19 与 2.7.9，这个相容族由同构
\begin{equation*}
\cO_{X\times_S Y}/\sJ_f\cong\cO_{X\times_S Y}/\sJ_g
\end{equation*}
诱导，使得下图交换：
\begin{center}
\begin{tikzcd}
\cO_{X\times_S Y} \arrow[r] \arrow[dr] & \cO_{X\times_S Y}/\sJ_f \arrow[d,"\cong"] \\
& \cO_{X\times_S Y}/\sJ_g.
\end{tikzcd}
\end{center}
所以我们必有
\begin{equation*}
\cO_{X\times_S Y}/\sJ_f=\cO_{X\times_S Y}/\sJ_g.
\end{equation*}
这说明 $f$ 与 $g$ 有相同的图像，从而 $f=g$。所以标准映射
\begin{equation*}
\Hom_S(X,Y)\to\plim_k\Hom_{S_k}(X_k,Y_k)
\end{equation*}
是单射。
设 $f_k\colon X_k\to Y_k$（$k\in\bN$）是一族 $S_k$-态射，使得对每对 $l\ge k$，$f_k$ 由 $f_l$ 经基变换 $S_k\to S_l$ 诱导。设
\begin{equation*}
\Gamma_{f_k}\colon X_k\to X_k\times_{S_k}Y_k
\end{equation*}
是 $f_k$ 的图像，它必是闭浸入，设 $\cJ_k$ 是其理想层。则族 $(\cO_{X_k\times_{S_k}Y_k}/\cJ_k)_{k\in\bN}$ 满足 2.7.11 的条件。于是存在凝聚 $\cO_{X\times_S Y}$-模 $\cF$，使得族 $(\cO_{X_k\times_{S_k}Y_k}/\cJ_k)_{k\in\bN}$ 同构于族 $(i_k^*\cF)_{k\in\bN}$。相容的满射族
\begin{equation*}
\cO_{X_k\times_{S_k}Y_k}\to\cO_{X_k\times_{S_k}Y_k}/\cJ_k
\end{equation*}
可以等同于相容的满射族
\begin{equation*}
i_k^*\cO_{X\times_S Y}\to i_k^*\cF.
\end{equation*}
由 1.5.19 与 2.7.9，这个满射族由满射
\begin{equation*}
\cO_{X\times_S Y}\to\cF
\end{equation*}
诱导。设 $\cJ$ 是其核，设 $Z$ 是 $X\times_S Y$ 的由理想层 $\cJ$ 定义的闭子概形。我们有同构
\begin{equation*}
i_k^*(\cO_{X\times_S Y}/\cJ)\cong i_k^*\cF\cong\cO_{X_k\times_{S_k}Y_k}/\cJ_k,
\end{equation*}
使得下图交换：
\begin{center}
\begin{tikzcd}
\cO_{X_k\times_{S_k}Y_k} \arrow[r,"\cong"] \arrow[d] & i_k^*\cO_{X\times_S Y} \arrow[r,"\cong"] \arrow[d] & i_k^*\cO_{X\times_S Y} \arrow[d] \\
\cO_{X_k\times_{S_k}Y_k}/\cJ_k \arrow[r,"\cong"] & i_k^*\cF \arrow[r,"\cong"] & i_k^*\cO_{X\times_S Y}/\cJ.
\end{tikzcd}
\end{center}
于是 $Z_k=Z\times_S S_k$ 是对应于 $f_k\colon X_k\to Y_k$ 的图像 $\Gamma_{f_k}$ 的闭子概形。设 $p_k\colon Z_k\to X_k$ 与 $p\colon Z\to X$ 是如下复合：
\begin{align*}
Z_k&\to X_k\times_{S_k}Y_k\to X_k,\\
Z&\to X\times_S Y\to X.
\end{align*}
由于 $Z_k$ 是 $f_k$ 的图像，每个 $p_k$ 是同构。由下面的 2.7.21，$p$ 也是同构。设 $f\colon X\to Y$ 是复合
\begin{equation*}
X\xrightarrow{p^{-1}}Z\to X\times_S Y\to Y.
\end{equation*}
则每个 $f_k\colon X_k\to Y_k$ 由 $f$ 经基变换 $S_k\to S$ 诱导。这就证明了映射
\begin{equation*}
\Hom_S(X,Y)\to\plim_k\Hom_{S_k}(X_k,Y_k)
\end{equation*}
是满射。
\end{proof}

\begin{lemma}\label{it:2.7.21}
设 $A$ 是诺特环，$I$ 是 $A$ 的理想，$S=\Spec A$，$Z$ 与 $X$ 是 $S$ 上的两个固有概形。对每个 $k$，设 $S_k=\Spec A/I^k$，$X_k=X\times_S S_k$，$Z_k=Z\times_S S_k$。设 $p\colon Z\to X$ 是 $S$-态射，使得对每个 $k$，由 $p$ 经基变换 $S_k\to S$ 诱导的态射 $p_k\colon Z_k\to X_k$ 是同构。若 $A$ 是 $I$-进完备的，则 $p$ 也是同构。
\end{lemma}

\begin{proof}
由 1.3.26 (vi)，$p\colon Z\to X$ 是固有的。设
\begin{equation*}
Z\xrightarrow{p'}X'\to X
\end{equation*}
是它的 Stein 分解。由于 $p_k\colon Z_k\to X_k$ 是同构，$Z$ 的任一落在 $Z_k$ 中的点 $z$ 在 $p^{-1}(p(z))$ 中是孤立的。由 2.5.13，存在 $Z_k$ 在 $Z$ 中的邻域 $Z'$，使得 $p'|_{Z'}\colon Z'\to X'$ 是开浸入。如同 2.7.9 的证明，$Z$ 是 $Z_k$ 在 $Z$ 中的唯一邻域。于是 $p'\colon Z\to X'$ 是开浸入。但 $p'$ 是固有且满射的。所以 $p'$ 是同构。从而 $p\colon Z\to X$ 是有限态射。由 1.4.11，要证明 $p$ 是同构，只需证明
\begin{equation*}
p^{\sharp}\colon\cO_X\to p_*\cO_Z
\end{equation*}
是同构。由 2.7.9、1.5.17 与 1.5.19，只需证明对每个 $k$，
\begin{equation*}
i_k^*\cO_X\to i_k^*p_*\cO_Z
\end{equation*}
是同构，其中 $i_k\colon X_k\to X$ 是闭浸入。利用 $p$ 是仿射态射这一事实，可以验证
\begin{equation*}
i_k^*p_*\cO_Z\cong p_{k*}\cO_{Z_k}.
\end{equation*}
于是态射
\begin{equation*}
i_k^*\cO_X\to i_k^*p_*\cO_Z
\end{equation*}
可以等同于
\begin{equation*}
p_k^{\sharp}\colon\cO_{X_k}\to p_{k*}\cO_{Z_k}.
\end{equation*}
由于 $p_k$ 是同构，$p_k^{\sharp}$ 也是。这就证明了我们的断言。
\end{proof}

\begin{tnote}
原书 p.245 在 2.7.21 的结论里印作「then $f$ is also an isomorphism」，但该引理通篇只有 $p\colon Z\to X$，没有 $f$，其证明也是从 $p$ 出发。译文按 $p$ 处理。[依据：原书 p.245 引理 2.7.21 及其证明]
\end{tnote}

\begin{tnote}
至此全书正文结束：原书 p.246 在 2.7.21 的证明后即止，§2.7 没有配习题。[依据：原书 p.246]
\end{tnote}
